%%%%%%%%%%%%%%%%%%%%%%%%%%%%%%%%%%%%%%%%%%%%%%%%%%%%%%%%%%%%%%%%%%%%%%%%%%%%%%%
% SECTION 12.2: The Thinking Style Parameter φ
% Module: BCM2_06_WAX_7250
% EBF Question 1: HOW does he/she think?
%%%%%%%%%%%%%%%%%%%%%%%%%%%%%%%%%%%%%%%%%%%%%%%%%%%%%%%%%%%%%%%%%%%%%%%%%%%%%%%

\section{Question 1: How Does the Person Think? ($\varphi$)}
\label{sec:12-thinking-style}

The first question this chapter answers is fundamental: When facing a decision 
with multiple dimensions — financial, social, environmental, identity — 
\textbf{how does the person integrate them?}

%------------------------------------------------------------------------------
\subsection{The Core Insight: Two Ways of Thinking}
\label{subsec:12-2-core-insight}
%------------------------------------------------------------------------------

\begin{intuition}
\textbf{Consider two people buying a heat pump:}

\textbf{Maria} ($\varphi \to 1$): ``I'll only buy if the price is right AND 
it's good for the environment AND my neighbors approve AND it fits my identity 
as a responsible person. If ANY of these fails, I won't do it.''

\textbf{Thomas} ($\varphi \to 0$): ``I look at the total picture. Great financial 
savings can make up for the fact that my neighbors haven't gotten one yet. 
I add up all the pros and cons.''

These represent fundamentally different \textit{thinking styles} — and they 
lead to different decisions even with identical underlying utilities.
\end{intuition}

%------------------------------------------------------------------------------
\subsection{Definition: The Complementarity Parameter $\varphi$}
\label{subsec:12-2-definition}
%------------------------------------------------------------------------------

\begin{definition}[Complementarity Parameter]
\label{def:phi}
The \textbf{Complementarity Parameter} $\varphi \in [0,1]$ captures how a person 
integrates multiple utility dimensions when making a decision:
\begin{itemize}
\item $\varphi = 1$: Perfect complementarity — all dimensions must be satisfied
\item $\varphi = 0$: Perfect substitutability — dimensions compensate each other
\item $\varphi \in (0,1)$: Mixed — partial complementarity
\end{itemize}
\end{definition}

\begin{equation}
\boxed{\varphi \in [0, 1]}
\end{equation}

\begin{table}[htbp]
\centering
\caption{The Two Thinking Styles}
\label{tab:thinking-styles}
\begin{tabular}{lp{5.5cm}p{5.5cm}}
\toprule
& \textbf{Connected Thinking} ($\varphi \to 1$) & \textbf{Separated Thinking} ($\varphi \to 0$) \\
\midrule
\textbf{Logic} & Complementary (AND) & Additive (SUM) \\[4pt]
\textbf{Metaphor} & Chain: only as strong as weakest link & Portfolio: diversification works \\[4pt]
\textbf{Formula} & $\min\{U_1^*, U_2^*, U_3^*\}$ & $U_1^* + U_2^* + U_3^*$ \\[4pt]
\textbf{Weak dimension} & Blocks everything & Gets compensated \\[4pt]
\textbf{Example} & ``ALL must be right'' & ``Total benefit counts'' \\[4pt]
\textbf{Risk profile} & Risk-averse, thorough & Trade-off accepting \\
\bottomrule
\end{tabular}
\end{table}

%------------------------------------------------------------------------------
\subsection{Why $\varphi$ Matters: A Numerical Example}
\label{subsec:12-2-why-matters}
%------------------------------------------------------------------------------

\begin{example}[Same Utilities, Different Decisions]
Consider two people with identical effective utilities:
\begin{center}
\begin{tabular}{lccc}
\toprule
& $U_{INU}^*$ & $U_{KNU}^*$ & $U_{IDU}^*$ \\
\midrule
Both persons & 0.8 & 0.7 & 0.3 \\
\bottomrule
\end{tabular}
\end{center}

Now calculate WAX for different $\varphi$ values:

\textbf{Maria} ($\varphi = 0.8$, connected thinking):
\begin{align*}
WAX &= 0.8 \cdot \min\{0.8, 0.7, 0.3\} + 0.2 \cdot (0.8 + 0.7 + 0.3)\\
&= 0.8 \cdot 0.3 + 0.2 \cdot 1.8\\
&= 0.24 + 0.36 = \mathbf{0.60}
\end{align*}

\textbf{Thomas} ($\varphi = 0.2$, separated thinking):
\begin{align*}
WAX &= 0.2 \cdot \min\{0.8, 0.7, 0.3\} + 0.8 \cdot (0.8 + 0.7 + 0.3)\\
&= 0.2 \cdot 0.3 + 0.8 \cdot 1.8\\
&= 0.06 + 1.44 = \mathbf{1.50}
\end{align*}

\textbf{Result:} With identical utilities, Thomas has 2.5× higher WAX than Maria!

If $\theta = 1.0$:
\begin{itemize}
\item Maria: $WAX = 0.60 < \theta$ → Won't act
\item Thomas: $WAX = 1.50 > \theta$ → Will act
\end{itemize}
\end{example}

\begin{intuition}
\textbf{The Key Insight:}

For connected thinkers ($\varphi \to 1$), the \textbf{weakest dimension} 
dominates. Maria's low identity utility (0.3) drags down her entire decision.

For separated thinkers ($\varphi \to 0$), \textbf{strong dimensions compensate} 
for weak ones. Thomas's excellent financial utility (0.8) outweighs his 
identity concerns.

\textit{Intervention implication:} For Maria, you must fix the weakest link. 
For Thomas, you can strengthen any dimension.
\end{intuition}

%------------------------------------------------------------------------------
\subsection{$\varphi$ Is Not Fixed: The Decomposition}
\label{subsec:12-2-decomposition}
%------------------------------------------------------------------------------

A person's thinking style is not a fixed trait. It shifts with context:

\begin{equation}
\boxed{\varphi = \varphi_0 + \Delta\varphi(\Psi) = \varphi_0 + \sum_j \gamma_j \cdot \Psi_j}
\label{eq:phi-decomposition}
\end{equation}

where:
\begin{itemize}
\item $\varphi_0$ = \textbf{Baseline} thinking style (personality, relatively stable)
\item $\Delta\varphi(\Psi)$ = \textbf{Context-induced shift}
\item $\gamma_j$ = Sensitivity of $\varphi$ to context dimension $\Psi_j$
\end{itemize}

\begin{intuition}
\textbf{What shifts $\varphi$?}

Your thinking style changes depending on the situation:

\begin{center}
\begin{tabular}{lcp{6cm}}
\toprule
\textbf{Context} & \textbf{Effect on $\varphi$} & \textbf{Why?}\\
\midrule
High uncertainty & $\downarrow$ & Under stress, fall back to self-interest\\
Strong group identity & $\uparrow$ & ``What would my group think?'' — norms bind\\
Identity-salient decision & $\uparrow$ & ``Who am I?'' becomes central\\
Being observed & $\uparrow$ & Social dimensions become non-negotiable\\
Time pressure & $\downarrow$ & No time for connected evaluation\\
\bottomrule
\end{tabular}
\end{center}
\end{intuition}

%------------------------------------------------------------------------------
\subsection{The $\gamma$-Matrix: Quantifying Context Effects}
\label{subsec:12-2-gamma-matrix}
%------------------------------------------------------------------------------

The $\gamma_j$ coefficients quantify how each context dimension shifts $\varphi$:

\begin{table}[htbp]
\centering
\caption{The $\gamma$-Matrix: Context Effects on Thinking Style}
\label{tab:gamma-matrix}
\begin{tabular}{lccl}
\toprule
\textbf{Context $\Psi_j$} & \textbf{$\gamma_j$} & \textbf{Direction} & \textbf{Mechanism}\\
\midrule
$\Psi_{UNS}$ (Uncertainty) & $-0.15$ & $\downarrow$ & Retreat to self-interest\\
$\Psi_{KOM}$ (Complexity) & $-0.10$ & $\downarrow$ & Cognitive overload → simplify\\
$\Psi_{SEG}$ (Segmentation) & $+0.20$ & $\uparrow$ & Group norms become binding\\
$\Psi_{IDN}$ (Identity) & $+0.25$ & $\uparrow$ & Self-concept central\\
$\Psi_{SOC}$ (Social) & $+0.15$ & $\uparrow$ & Others are watching\\
$\Psi_{REG}$ (Regulation) & $+0.05$ & $\uparrow$ & External structure helps\\
$\Psi_{TMP}$ (Temporal) & $-0.08$ & $\downarrow$ & Time pressure → shortcuts\\
\bottomrule
\end{tabular}
\end{table}

\begin{example}[Anna's $\varphi$-Shift]
Anna has baseline $\varphi_0 = 0.35$ (naturally additive thinker).

\textbf{Situation:} Her neighbor mentions ``Everyone on our street is getting 
heat pumps.'' This activates:
\begin{itemize}
\item $\Psi_{SEG} = 0.8$ (strong neighborhood identity)
\item $\Psi_{SOC} = 0.7$ (social observation)
\end{itemize}

\textbf{Calculation:}
\begin{align*}
\Delta\varphi &= \gamma_{SEG} \cdot \Psi_{SEG} + \gamma_{SOC} \cdot \Psi_{SOC}\\
&= 0.20 \cdot 0.8 + 0.15 \cdot 0.7\\
&= 0.16 + 0.105 = 0.265
\end{align*}

\textbf{Result:}
\begin{equation*}
\varphi_{new} = 0.35 + 0.265 = \mathbf{0.615}
\end{equation*}

Anna shifted from separated thinking (0.35) toward connected thinking (0.615). 
Now the social dimension becomes more binding — she can't just ignore what 
the neighbors think.
\end{example}

%------------------------------------------------------------------------------
\subsection{Estimating $\varphi$: The LLM-MC Method}
\label{subsec:12-2-estimation}
%------------------------------------------------------------------------------

How do we measure $\varphi$ for a specific person in a specific context?

\begin{note}[LLM-MC Estimation]
The \textbf{Large Language Model Monte Carlo} (LLM-MC) method enables 
$\varphi$-estimation through:
\begin{enumerate}
\item Constructing a detailed persona (demographics, values, context)
\item Presenting multi-dimensional trade-off scenarios
\item Running 1000+ simulated decisions with temperature variation
\item Inferring $\varphi$ from the pattern of choices
\end{enumerate}
\end{note}

\textbf{Formal specification:} See Appendix AV, Section ``$\varphi$-Estimation 
via LLM-MC'' for the complete methodology, validation against experimental 
data, and calibration procedures.

%------------------------------------------------------------------------------
\subsection{Summary: Question 1 Answered}
\label{subsec:12-2-summary}
%------------------------------------------------------------------------------

\begin{tcolorbox}[
  colback=blue!3,
  colframe=blue!50!black,
  title={\textbf{Question 1: How Does the Person Think?}},
  fonttitle=\bfseries\sffamily
]

\textbf{The Answer:} The Complementarity Parameter $\varphi$

\begin{equation*}
\varphi = \varphi_0 + \sum_j \gamma_j \cdot \Psi_j
\end{equation*}

\begin{tabular}{lp{9cm}}
$\varphi \to 1$ & Connected: ``Everything must be right'' — weakest link dominates\\
$\varphi \to 0$ & Separated: ``Total counts'' — dimensions compensate\\
\end{tabular}

\vspace{6pt}
\textbf{Key insight:} The same utilities lead to different decisions 
depending on thinking style.

\textbf{Intervention lever:} Shift $\varphi$ by changing context 
(social norms, identity salience, uncertainty reduction).

\end{tcolorbox}

\begin{cross-reference}
\textbf{Formal Specification:} See Appendix AV for complete $\gamma$-matrix, 
LLM-MC estimation methodology, and empirical validation against 
Henrich et al. (2010) cross-cultural data.
\end{cross-reference}

